%%------------------------------------------------------------------------------
%%Präambel
%%------------------------------------------------------------------------------
%%Klasse
%%------------------------------------------------------------------------------
\documentclass[a5paper,landscape,DIV=calc,11pt,notitlepage]{scrartcl}
%%------------------------------------------------------------------------------
%%Packages
%%------------------------------------------------------------------------------
\usepackage[headsepline=no, footsepline=no]{scrlayer-scrpage}
\usepackage[T1]{fontenc}
\usepackage{libertinus}
\usepackage[ngerman]{babel} 
\usepackage[right=24pt,left=24pt, top=24pt, bottom=24pt,twoside=false]{geometry}
\usepackage{tabularx}
\usepackage[svgnames,RGB,table]{xcolor}
\usepackage{microtype}
\usepackage[german]{layout}
\usepackage{tcolorbox}
\usepackage{hyperref}
\hypersetup{colorlinks, citecolor=, linkcolor=, urlcolor=blue,pdfauthor={Sascha Christmann},pdftitle={G\={o}j\={u}-Ry\={u} Karate-D\={o} Reusrath - Prüfungsordnung - Auszug},pdfsubject={G\={o}j\={u}-Ry\={u} Karate-D\={o}},pdfproducer={MiKTeX},pdfcreator={pdflatex}}
\usepackage{enumitem}
\usepackage{tabto}
%%------------------------------------------------------------------------------
%%Vereinbarungen & Makros
%%------------------------------------------------------------------------------
\clearpairofpagestyles
%%------------------------------------------------------------------------------
\definecolor{YBELT}{rgb}{1,1,0.2}
\definecolor{OBELT}{rgb}{1,0.404,0}
\definecolor{GBELT}{rgb}{0.431,0.686,0.282}
\definecolor{BLBELT}{rgb}{0,0.408,1}
\definecolor{BRBELT}{rgb}{0.376,0.22,0.075}
%%------------------------------------------------------------------------------
\tcbset{width=\textwidth,height=\textheight,right=12pt,left=12pt,colback=white,fonttitle=\bfseries,adjust text=Äpgjy}
%%------------------------------------------------------------------------------
%%Ende Präambel
%%------------------------------------------------------------------------------
\begin{document}
%%------------------------------------------------------------------------------
%%Karte_Gelb
%%------------------------------------------------------------------------------
\begin{tcolorbox}[colback=white,colframe=YBELT,coltitle=black,title=8. Kyu: Gelb]
	\renewcommand{\labelenumi}{\alph{enumi}.)}
	\renewcommand{\labelenumii}{\arabic*.}
	{\scriptsize\textbf{\textit{Der 8.\,Kyu wird in aller Regel gemeinsam mit dem 9.\,Kyu absolviert, daher ist für diese Prüfung die doppelte Prüfungsgebühr, für entsprechend 2 Prüfungsmarken, zu entrichten.}}\newline Anmerkung zum \textit{(Chisai-no-)}Mawashi-Geri: in der Chisai-Form wird diese Technik als kleiner Halbkreisfußtritt ohne Drehung des Standbeins ausgeführt und ist technisch vollkommen ausreichend. Die sportliche Variante kann und darf aber auch trainiert und gezeigt werden.\newline
	\begin{tabularx}{\textwidth}{lX}
		\textit{Vorausgesetzte Techniken}:& Zuki, Age-Uke, Yoko-Uke, Harai-Otoshi-Uke, Soto-Uke, Mae-Geri, \textit{(Chisai-no-)}Mawashi-Geri
	\end{tabularx}\\}
	\begin{enumerate}[nosep]
		\item \textit{\textbf{Kihon-Ido}}:
		\begin{enumerate}[nosep]
			\item Harai-Otoshi-Uke/Gyaku-Zuki Ch\={u}dan\tabto{280pt}(Zenkutsu-Dachi)
			\item Mae-Geri Ch\={u}dan\tabto{280pt}(Zenkutsu-Dachi)
			\item \textit{(Chisai-no-)}Mawashi-Geri Ch\={u}dan\tabto{280pt}(Zenkutsu-Dachi)
			\item Oi-Zuki J\={o}dan/Gyaku-Zuki Ch\={u}dan\tabto{280pt}(Zenkutsu-Dachi)
			\item Soto-Uke Ch\={u}dan/Gyaku-Zuki Ch\={u}dan\tabto{280pt}(Zenkutsu-Dachi)
			\item Age-Uke/Gyaku-Zuki J\={o}dan\tabto{280pt}(Sanchin-Dachi)
			\item Yoko-Uke/Gyaku-Zuki Ch\={u}dan\tabto{280pt}(Sanchin-Dachi)
		\end{enumerate}
		\item \textit{\textbf{Kata}}:
		\begin{enumerate}[nosep]
			\item Taikyoku J\={o}dan
			\item Taikyoku Ch\={u}dan
		\end{enumerate}
		\item \textit{\textbf{Partnerformen}}:
		\begin{enumerate}[nosep]
			\item Kihon-Ippon Kumite:\tabto{280pt}J\={o}dan\&Ch\={u}dan (ggf. Gedan)
		\end{enumerate}
	\end{enumerate}
\end{tcolorbox}
%%------------------------------------------------------------------------------
\pagebreak
%%------------------------------------------------------------------------------
%%Karte_Orange
%%------------------------------------------------------------------------------
\begin{tcolorbox}[colback=white,colframe=OBELT,coltitle=white,title=7. Kyu: Orange]
	\renewcommand{\labelenumi}{\alph{enumi}.)}
	\renewcommand{\labelenumii}{\arabic*.}
	{\scriptsize
		\begin{tabularx}{\textwidth}{lX}
			\textit{Vorausgesetzte Techniken}:& Zuki, Age-Uke, Yoko-Uke, Harai-Otoshi-Uke, Soto-Uke, Uraken-Uchi, Empi-Age-Uchi, Mae-Geri, \mbox{\textit{(Chisai-no-)}Mawashi-Geri}
		\end{tabularx}\\}
	\begin{enumerate}[nosep]
		\item \textit{\textbf{Kihon-Ido}}:
		\begin{enumerate}[nosep]
			\item Harai-Otoshi-Uke/Gyaku-Zuki Ch\={u}dan\tabto{280pt}(Zenkutsu-Dachi)
			\item Oi-Zuki J\={o}dan/Gyaku-Zuki Ch\={u}dan\tabto{280pt}(Zenkutsu-Dachi)
			\item Soto-Uke Ch\={u}dan/Gyaku-Zuki Ch\={u}dan\tabto{280pt}(Zenkutsu-Dachi)
			\item Kansetsu-Geri Gedan\tabto{280pt}(Sanchin-Dachi)
			\item Age-Uke/Gyaku-Zuki J\={o}dan\tabto{280pt}(Sanchin-Dachi)
			\item Yoko-Uke/Gyaku-Zuki Ch\={u}dan\tabto{280pt}(Sanchin-Dachi)
			\item Harai-Otoshi-Uke/Gyaku-Zuki Gedan\tabto{280pt}(Shiko-Dachi 45°)
		\end{enumerate}
		\item \textit{\textbf{Kata}}:
		\begin{enumerate}[nosep]
			\item Taikyoku (welche Taikyoku geprüft wird, entscheidet der Prüfer)
			\item Gekisai-Dai-Ichi
		\end{enumerate}
		\item \textit{\textbf{Partnerformen}}:
		\begin{enumerate}[nosep]
			\item Kihon-Ippon Kumite:\tabto{280pt}J\={o}dan\&Ch\={u}dan (ggf. Gedan)
		\end{enumerate}
	\end{enumerate}
\end{tcolorbox}
%%------------------------------------------------------------------------------
\pagebreak
%%------------------------------------------------------------------------------
%%Karte_Grün
%%------------------------------------------------------------------------------
\begin{tcolorbox}[colback=white,colframe=GBELT,coltitle=white,title=6. Kyu: Grün]
	\renewcommand{\labelenumi}{\alph{enumi}.)}
	\renewcommand{\labelenumii}{\arabic*.}
	{\scriptsize
		\begin{tabularx}{\textwidth}{lX}
			\textit{Vorausgesetzte Techniken}:& Zuki, Age-Uke, Yoko-Uke, Harai-Otoshi-Uke, Soto-Uke,Kake-Uke, Mawashi-Uke, Uraken-Uchi, Empi-Age-Uchi, Mae-Geri, \textit{(Chisai-no-)}Mawashi-Geri
		\end{tabularx}\\}
	\begin{enumerate}[nosep]
		\item \textit{\textbf{Kihon-Ido}}:
		\begin{enumerate}[nosep]
			\item Harai-Otoshi-Uke/Mae-Geri Ch\={u}dan/Gyaku-Zuki Ch\={u}dan\tabto{280pt}(Zenkutsu-Dachi)
			\item \textit{(Chisai-no-)}Mawashi-Geri Ch\={u}dan/Gyaku-Zuki Ch\={u}dan\tabto{280pt}(Zenkutsu-Dachi)
			\item Age-Uke/Gyaku-Zuki Ch\={u}dan/Mae-Geri Ch\={u}dan\tabto{280pt}(Zenkutsu-Dachi)
			\item Yoko-Uke/Ren-Zuki Jodon-Ch\={u}dan\tabto{280pt}(Sanchin-Dachi)
			\item Harai-Otoshi-Uke/Gyaku-Zuki Gedan\tabto{280pt}(Shiko-Dachi 45°)
			\item Kake-Uke/Mae-Geri Ch\={u}dan\tabto{280pt}(Neko-Ashi-Dachi)
		\end{enumerate}
		\item \textit{\textbf{Kata}}:
		\begin{enumerate}[nosep]
			\item Gekisai-Dai-Ichi
			\item Gekisai-Dai-Ni
		\end{enumerate}
		\item \textit{\textbf{Partnerformen (2 aus 3)}}:
		\begin{enumerate}[nosep]
			\item Yakusoku-Kumite:\tabto{200pt}3 Kumite-Ura\,\&\,3 Nage-Waza (frei auszuwählen)
			\item Shiai-Kumite:\tabto{200pt}Jiyu-Kumite: 1 Minute \frqq Freikampf\flqq~(sportlich)
			\item Alternativ zum Shiai-Kumite:\tabto{200pt}SV
		\end{enumerate}
	\end{enumerate}
\end{tcolorbox}
%%------------------------------------------------------------------------------
\pagebreak
%%------------------------------------------------------------------------------
%%Karte_Blau_1
%%------------------------------------------------------------------------------
\begin{tcolorbox}[colback=white,colframe=BLBELT,coltitle=white,title=5. Kyu: Blau]
	\renewcommand{\labelenumi}{\alph{enumi}.)}
	\renewcommand{\labelenumii}{\arabic*.}
	{\scriptsize
		\begin{tabularx}{\textwidth}{lX}
			\textit{Vorausgesetzte Techniken}:& Zuki, Age-Uke, Yoko-Uke, Harai-Otoshi-Uke, Soto-Uke, Kake-Uke, Mawashi-Uke, Teisho-Uchi, Uraken-Uchi, \mbox{Empi-Age-Uchi}, Mae-Geri, \textit{(Chisai-no-)}Mawashi-Geri
		\end{tabularx}\\}
	\begin{enumerate}[nosep]
		\item \textit{\textbf{Kihon-Ido}}:
		\begin{enumerate}[nosep]
			\item Harai-Otoshi-Uke/Mae-Geri Ch\={u}dan/Gyaku-Zuki Ch\={u}dan\tabto{280pt}(Zenkutsu-Dachi)
			\item \textit{(Chisai-no-)}Mawashi-Geri Ch\={u}dan/Gyaku-Zuki Ch\={u}dan\tabto{280pt}(Zenkutsu-Dachi)
			\item Age-Uke/Gyaku-Zuki Ch\={u}dan/Mae-Geri Ch\={u}dan\tabto{280pt}(Zenkutsu-Dachi)
			\item Yoko-Uke/Ren-Zuki Jodon-Ch\={u}dan\tabto{280pt}(Sanchin-Dachi)
			\item Harai-Otoshi-Uke/Gyaku-Zuki Gedan\tabto{280pt}(Shiko-Dachi 45°)
			\item Kake-Uke/Mae-Geri Ch\={u}dan\tabto{280pt}(Neko-Ashi-Dachi)
		\end{enumerate}
		\item \textit{\textbf{Kata}}:
		\begin{enumerate}[nosep]
			\item Gekisai-Dai-Ni
			\item Saifa
		\end{enumerate}
		\item \textit{\textbf{Partnerformen (2 aus 3)}}:
		\begin{enumerate}[nosep]
			\item Yakusoku-Kumite:\tabto{200pt}3 Kumite-Ura\,\&\,3 Nage-Waza (frei auszuwählen)
			\item Shiai-Kumite:\tabto{200pt}Jiyu-Kumite: 1 Minute \frqq Freikampf\flqq~(sportlich)
			\item Alternativ zum Shiai-Kumite:\tabto{200pt}SV
		\end{enumerate}
	\end{enumerate}
\end{tcolorbox}
%%------------------------------------------------------------------------------
\pagebreak
%%------------------------------------------------------------------------------
%%Karte_Blau_2
%%------------------------------------------------------------------------------
\begin{tcolorbox}[colback=white,colframe=BLBELT,coltitle=white,title=4. Kyu: Blau]
	\renewcommand{\labelenumi}{\alph{enumi}.)}
	\renewcommand{\labelenumii}{\arabic*.}
	{\scriptsize
		\begin{tabularx}{\textwidth}{lX}
			\textit{Vorausgesetzte Techniken}:& Zuki, Age-Uke, Yoko-Uke, Harai-Otoshi-Uke, Soto-Uke, Kake-Uke, Mawashi-Uke, Teisho-Uchi, Uraken-Uchi, \mbox{Empi-Age-Uchi}, Mae-Geri, \textit{(Chisai-no-)}Mawashi-Geri
		\end{tabularx}\\}
	\begin{enumerate}[nosep]
		\item \textit{\textbf{Kihon-Ido}}:
		\begin{enumerate}[nosep]
			\item Harai-Otoshi-Uke/Mae-Geri Ch\={u}dan/Gyaku-Zuki Ch\={u}dan\tabto{280pt}(Zenkutsu-Dachi)
			\item \textit{(Chisai-no-)}Mawashi-Geri Ch\={u}dan/Gyaku-Zuki Ch\={u}dan\tabto{280pt}(Zenkutsu-Dachi)
			\item Age-Uke/Gyaku-Zuki Ch\={u}dan/Mae-Geri Ch\={u}dan\tabto{280pt}(Zenkutsu-Dachi)
			\item Yoko-Uke/Ren-Zuki Jodon-Ch\={u}dan\tabto{280pt}(Sanchin-Dachi)
			\item Harai-Otoshi-Uke/Gyaku-Zuki Gedan\tabto{280pt}(Shiko-Dachi 45°)
			\item Kake-Uke/Mae-Geri Ch\={u}dan\tabto{280pt}(Neko-Ashi-Dachi)
		\end{enumerate}
		\item \textit{\textbf{Kata}}:
		\begin{enumerate}[nosep]
			\item Saifa
			\item Seenchin
		\end{enumerate}
		\item \textit{\textbf{Partnerformen (2 aus 3)}}:
		\begin{enumerate}[nosep]
			\item Yakusoku-Kumite:\tabto{200pt}3 Kumite-Ura\,\&\,3 Nage-Waza (frei auszuwählen)
			\item Shiai-Kumite:\tabto{200pt}Jiyu-Kumite: 1 Minute \frqq Freikampf\flqq~(sportlich)
			\item Alternativ zum Shiai-Kumite:\tabto{200pt}SV
		\end{enumerate}
	\end{enumerate}
\end{tcolorbox}
%%------------------------------------------------------------------------------
\pagebreak
%%------------------------------------------------------------------------------
%%Karte_Braun_1
%%------------------------------------------------------------------------------
\begin{tcolorbox}[colback=white,colframe=BRBELT,coltitle=white,title=3. Kyu: Braun]
	\renewcommand{\labelenumi}{\alph{enumi}.)}
	\renewcommand{\labelenumii}{\arabic*.}
	{\scriptsize 
		\begin{tabularx}{\textwidth}{lX}
			\textit{Vorausgesetzte Techniken}:& Zuki, Age-Uke, Yoko-Uke, Harai-Otoshi-Uke, Soto-Uke, Kake-Uke, Mawashi-Uke, Teisho-Uchi, Uraken-Uchi, \mbox{Empi-Age-Uchi}, Mae-Geri, \textit{(Chisai-no-)}Mawashi-Geri
		\end{tabularx}\\}\vspace{12pt}
	\begin{enumerate}[nosep]
		\item \textit{\textbf{Kihon-Ido}}:
		\begin{enumerate}[nosep]
			\item Harai-Otoshi-Uke/Mae-Geri Ch\={u}dan/Gyaku-Zuki Ch\={u}dan\tabto{280pt}(Zenkutsu-Dachi)
			\item J\={o}dan Kizami-Zuki/Ch\={u}dan Gyku-Zuki\tabto{280pt}(Suri-Ashi/Zenkutsu-Dachi)
			\item Ch\={u}dan Soto-Uke/J\={o}dan Uraken-Uchi\tabto{280pt}(Sanchin-Dachi)
			\item Kake-Uke/Ren-Zuki J\={o}dan-Ch\={u}dan\tabto{280pt}(Sanchin-Dachi)
			\item Harai-Otoshi-Uke/Gyaku-Zuki Gedan\tabto{280pt}(Shiko-Dachi 45°)
			\item Shuto-Uke/Gedan Kansetsu-Geri\tabto{280pt}(Neko-Ashi-Dachi)
		\end{enumerate}\vspace{24pt}
		\item \textit{\textbf{Kata}}:\tabto{200pt}\textit{\textbf{Kata\,Bunkai}}:
		\begin{enumerate}[nosep]
			\item Seenchin\tabto{200pt}Geki-Sai-Dai-Ichi
			\item Tensho
		\end{enumerate}\vspace{24pt}
		\item \textit{\textbf{Partnerformen (2 aus 3)}}:
		\begin{enumerate}[nosep]
			\item Yakusoku-Kumite:\tabto{200pt}3 Kumite-Ura\,\&\,3 Nage-Waza (frei auszuwählen)
			\item Shiai-Kumite:\tabto{200pt}Jiyu-Kumite: 1 Minute \frqq Freikampf\flqq~(sportlich)
			\item Goshin-Jitsu-Kumite:\tabto{200pt}SV gegen Halten, Würgen, Stoßen
		\end{enumerate}
	\end{enumerate}
\end{tcolorbox}
%%------------------------------------------------------------------------------
\pagebreak
%%------------------------------------------------------------------------------
%%Karte_Braun_2
%%------------------------------------------------------------------------------
\begin{tcolorbox}[colback=white,colframe=BRBELT,coltitle=white,title=2. Kyu: Braun]
	\renewcommand{\labelenumi}{\alph{enumi}.)}
	\renewcommand{\labelenumii}{\arabic*.}
	{\scriptsize 
		\begin{tabularx}{\textwidth}{lX}
			\textit{Vorausgesetzte Techniken}:& Zuki, Age-Uke, Yoko-Uke, Harai-Otoshi-Uke, Soto-Uke, Kake-Uke, Mawashi-Uke, Teisho-Uchi, Uraken-Uchi, \mbox{Empi-Age-Uchi}, Mae-Geri, \textit{(Chisai-no-)}Mawashi-Geri
		\end{tabularx}\\}\vspace{12pt}
	\begin{enumerate}[nosep]
		\item \textit{\textbf{Kihon-Ido}}:
		\begin{enumerate}[nosep]
			\item Harai-Otoshi-Uke/Mae-Geri Ch\={u}dan/Gyaku-Zuki Ch\={u}dan\tabto{280pt}(Zenkutsu-Dachi)
			\item J\={o}dan Kizami-Zuki/Ch\={u}dan Gyku-Zuki\tabto{280pt}(Suri-Ashi/Zenkutsu-Dachi)
			\item Ch\={u}dan Soto-Uke/J\={o}dan Uraken-Uchi\tabto{280pt}(Sanchin-Dachi)
			\item Kake-Uke/Ren-Zuki J\={o}dan-Ch\={u}dan\tabto{280pt}(Sanchin-Dachi)
			\item Harai-Otoshi-Uke/Gyaku-Zuki Gedan\tabto{280pt}(Shiko-Dachi 45°)
			\item Shuto-Uke/Gedan Kansetsu-Geri\tabto{280pt}(Neko-Ashi-Dachi)
		\end{enumerate}\vspace{24pt}
		\item \textit{\textbf{Kata}}:\tabto{200pt}\textit{\textbf{Kata\,Bunkai}}:
		\begin{enumerate}[nosep]
			\item Tensho\tabto{200pt}Geki-Sai-Dai-Ni
			\item Sanseru
		\end{enumerate}\vspace{24pt}
		\item \textit{\textbf{Partnerformen (2 aus 3)}}:
		\begin{enumerate}[nosep]
			\item Yakusoku-Kumite:\tabto{200pt}3 Kumite-Ura\,\&\,3 Nage-Waza (frei auszuwählen)
			\item Shiai-Kumite:\tabto{200pt}Jiyu-Kumite: 1 Minute \frqq Freikampf\flqq~(sportlich)
			\item Goshin-Jitsu-Kumite:\tabto{200pt}SV gegen Halten, Würgen, Stoßen
		\end{enumerate}
	\end{enumerate}
\end{tcolorbox}
%%------------------------------------------------------------------------------
\pagebreak
%%------------------------------------------------------------------------------
%%Karte_Braun_3
%%------------------------------------------------------------------------------
\begin{tcolorbox}[colback=white,colframe=BRBELT,coltitle=white,title=1. Kyu: Braun]
	\renewcommand{\labelenumi}{\alph{enumi}.)}
	\renewcommand{\labelenumii}{\arabic*.}
	{\scriptsize 
		\begin{tabularx}{\textwidth}{lX}
			\textit{Vorausgesetzte Techniken}:& Zuki, Age-Uke, Yoko-Uke, Harai-Otoshi-Uke, Soto-Uke, Kake-Uke, Mawashi-Uke, Teisho-Uchi, Uraken-Uchi, \mbox{Empi-Age-Uchi}, Mae-Geri, \textit{(Chisai-no-)}Mawashi-Geri
		\end{tabularx}\\\vspace{12pt}}
	\begin{enumerate}[nosep]
		\item \textit{\textbf{Kihon-Ido}}:
		\begin{enumerate}[nosep]
			\item Harai-Otoshi-Uke/Mae-Geri Ch\={u}dan/Gyaku-Zuki Ch\={u}dan\tabto{280pt}(Zenkutsu-Dachi)
			\item J\={o}dan Kizami-Zuki/Ch\={u}dan Gyku-Zuki\tabto{280pt}(Suri-Ashi/Zenkutsu-Dachi)
			\item Ch\={u}dan Soto-Uke/J\={o}dan Uraken-Uchi\tabto{280pt}(Sanchin-Dachi, gleicher Arm)
			\item Kake-Uke/Ren-Zuki J\={o}dan-Ch\={u}dan\tabto{280pt}(Sanchin-Dachi)
			\item Harai-Otoshi-Uke/Gyaku-Zuki Gedan\tabto{280pt}(Shiko-Dachi 45°)
			\item Shuto-Uke/Gedan Kansetsu-Geri\tabto{280pt}(Neko-Ashi-Dachi)
		\end{enumerate}\vspace{24pt}
		\item \textit{\textbf{Kata}}:\tabto{200pt}\textit{\textbf{Kata\,Bunkai}}:
		\begin{enumerate}[nosep]
			\item Sanseru\tabto{200pt}Geki-Sai-Dai-Ichi \textit{oder} Geki-Sai-Dai-Ni
			\item Sanchin
		\end{enumerate}\vspace{24pt}
		\item \textit{\textbf{Partnerformen (2 aus 3)}}:
		\begin{enumerate}[nosep]
			\item Yakusoku-Kumite:\tabto{200pt}3 Kumite-Ura\,\&\,3 Nage-Waza (frei auszuwählen)
			\item Shiai-Kumite:\tabto{200pt}Jiyu-Kumite: 1 Minute \frqq Freikampf\flqq~(sportlich)
			\item Goshin-Jitsu-Kumite:\tabto{200pt}SV gegen Halten, Würgen, Stoßen
		\end{enumerate}
	\end{enumerate}
\end{tcolorbox}
%%------------------------------------------------------------------------------
\pagebreak
%%------------------------------------------------------------------------------
%%Karte_Schwarz
%%------------------------------------------------------------------------------
\begin{tcolorbox}[colback=white,colframe=black,coltitle=white,adjusted title=1. Dan: Schwarz]
	\renewcommand{\labelenumi}{\alph{enumi}.)}
	\renewcommand{\labelenumii}{\arabic*.}
	{\scriptsize 
		\begin{tabularx}{\textwidth}{lX}
			\textit{Vorausgesetzte Techniken}:& Zuki-Waza, Uke-Waza, Geri-Waza und Dachi-Waza der bis hierhin erlernten Kihon-Ido, Kata, Kata-Bunkai und Yakusoku-Kumite
		\end{tabularx}\\}
	\begin{enumerate}[nosep]
		\item \textit{\textbf{Kihon-Ido}}:
		\begin{enumerate}[nosep]
			\item Harai-Otoshi-Uke/Gyaku-Zuki J\={o}dan/Mae-Geri Ch\={u}dan\tabto{335pt}(Zenkutsu-Dachi)
			\item Age-Uke/Ren-Zuki J\={o}dan-Ch\={u}dan\tabto{335pt}(Sanchin-Dachi)
			\item Mawashi-Empi-Uke/Uraken-Uchi/Harai-Otoshi-Uke/Gyaku-Zuki Ch\={u}dan\tabto{335pt}(Shiko-Dachi 45°)
			\item Shuto-Uke/Mae-Geri Ch\={u}dan\tabto{335pt}(Neko-Ashi-Dachi)
		\end{enumerate}
		\vspace{12pt}
		\item \textit{\textbf{Kata}}:\tabto{200pt}\textit{\textbf{Kata\,Bunkai}}:
		\begin{enumerate}[nosep]
			\item Sanseru\tabto{200pt}Sanseru
		\end{enumerate}
		\vspace{12pt}
		\item \textit{\textbf{Partnerformen (2 aus 3)}}:
		\begin{enumerate}[nosep]
			\item Yakusoku-Kumite:\tabto{200pt}3 Kumite-Ura\,\&\,2 Nage-Waza (frei auszuwählen)
			\item Shiai-Kumite:\tabto{200pt}Jiyu-Kumite: 1 Minute \frqq Freikampf\flqq~(sportlich)\\\tabto{200pt}\textit{entfällt für Sportler über 35 Jahre}
			\item Goshin-Jitsu-Kumite:\tabto{200pt}Abwehr und Konter gegen Packen, Würgen und Halten\\\textit{(Jissen-Jitsu, Selbstverteidigung)}\tabto{200pt} von vorn, von der Seite und von hinten
		\end{enumerate}
		\vspace{12pt}
		\item \textit{\textbf{Weitere Verbandsvorgaben}}:
		\begin{enumerate}[nosep]
			\item Vorlage eines Kampfrichterlehrgangsnachweis
			\item Besuch eines Dansha-Lehrgangs der Stilrichtung G\={o}j\={u}-Ry\={u} Karate-D\={o}
			\item Unsere Abteilung empfiehlt die Teilnahme an Lehrgängen auf Landes-, oder Bundesebene, um den Horizont des eigenen Karate beständig zu erweitern
			\item Weitere Details in der Prüfungsordnung
		\end{enumerate}
	\end{enumerate}
\end{tcolorbox}
%%------------------------------------------------------------------------------
\end{document}
%%------------------------------------------------------------------------------
%%EOF
%%------------------------------------------------------------------------------